\scanpage{120}
s'annulent pas simultanément, i.e. ils engendrent $D^{\otimes-2}$ : on trouve alors pour $\rho_0$ une section de $\mathbb P^1_k$ prenant la valeur $\infty$ aux pts où $\delta_B$ s'annule, \deleted{la valeur} i.e. $H,H'$ contenus dans un $\overline{\mathrm m}$ Borel \inserted{(mais $H,H'$ étales)}, et la valeur $0$ en les pts où $H$ ou $H'$ n'est pas étale. \unsure{L}

\uline{NB} Le seul cas sur un corps où $\rho_0$ n'est pas défini, est celui où $B=B_D$ est une sous-alg. de Borel, pour droite $D\subsetA^{\natural}$ \unsure{tot.} isotr., et $H$ (disons, à moins d'échanger les rôles) est $k\oplus D$, et enfin $H'$ est \deleted{de la forme} un \deleted{vectoriel}\inserted{\unsure{2-}plan de $B$} qui contient $k$ et est distinct de $H$ sans plus.

\deleted{Cette \unsure{situat|situat.}} Deux telles situations sont conjuguées loc -- i.e. ce type de situation est standard [correspond au choix d'un Borel et d'un tore maximal dedans, i.e. \uline{couple de Killing}.]
% archive sheet number: 119
